\newpage
\section{Объём}

\begin{definition}
    Пусть $V$ --- евклидово пространство размерности $n$.
    \emph{Объём} --- это отображение $\omega: \underbrace{V \times V \times \cdots \times V}_{n} \to \mathbb{R}$, такое что:
    \begin{enumerate}
        \item \textbf{Кососимметричность:}
        \[
            \omega(v_1, \dots, v_i + \lambda v_j, \dots, v_n)
            = \omega(v_1, \dots, v_n),
            \quad i \neq j,\ \lambda \in \mathbb{R}
        \]
        \item \textbf{Однородность:}
        \[
            \omega(v_1, \dots, \lambda v_i, \dots, v_n)
            = \lambda\, \omega(v_1, \dots, v_n) \quad \lambda \in \mathbb{R}
        \]
    \end{enumerate}
\end{definition}

\begin{example}
    Геометрический смысл свойства~1: если к одному вектору прибавить
    кратное другого (сдвинуть сторону параллелограмма), объём не меняется.

    \begin{center}
    \begin{tikzpicture}[>=stealth, thick]
        % Исходный параллелограмм
        \fill[cbBlueBack] (0,0) -- (3,0) -- (3.8,1.5) -- (0.8,1.5) -- cycle;
        \draw[->] (0,0) -- (3,0) node[below] {$v_1$};
        \draw[->] (0,0) -- (0.8,1.5) node[above left] {$v_2$};
        \draw[dashed, cbGray] (3,0) -- (3.8,1.5) -- (0.8,1.5);
        \draw[<->] (4.1,0) -- (4.1,1.5) node[midway, right] {$h$};
        \node at (1.5, -0.6) {$S = v_1 \cdot h$};


        % Сдвинутый параллелограмм
        \fill[cbRoseBack, opacity=0.5] (6,0) -- (9,0) -- (11,1.5) -- (8,1.5) -- cycle;
        \draw[->] (6,0) -- (9,0) node[below] {$v_1$};
        \draw[->] (6,0) -- (8,1.5) node[above] {$v_2 + \lambda v_1$};
        \draw[dashed, cbGray] (9,0) -- (11,1.5) -- (8,1.5);
        \draw[<->] (11.2,0) -- (11.2,1.5) node[midway, right] {$h$};
        \node at (8.5, -0.6) {$S = v_1 \cdot h$};
    \end{tikzpicture}
    \end{center}

    Высота $h$ не меняется при сдвиге, поэтому площадь (и объём) одинаковы.
\end{example}



\begin{stmt}
    Пусть $\omega$ --- объём. Тогда:
    \begin{enumerate}[label=\textbf{\arabic*.}]
        \item $\omega$ полилинейна,
        \item $\omega$ знакопеременна:
        \[
            \omega(\dots, v_i, \dots, v_j, \dots) = -\omega(\dots, v_j, \dots, v_i, \dots),
        \]
        \item $\omega(v_1, \dots, v_n) = 0$, если $v_1, \dots, v_n$ линейно зависимы (\textcircled{$\Rightarrow$} выполняется тольео если $\omega \ \not\equiv \ 0$).
    \end{enumerate}

    \begin{proof}
        \textbf{\emph{3}} в \textcircled{$\Leftarrow$}: Пусть $v_1, \dots, v_n$ линейно зависимы. НУО $v_1 = \sum_{i=2}^{n} \alpha_i v_i$. Тогда:
    
        \[  
            \omega(v_1, \dots, v_n)
            = \omega\!\left(\sum_{i=2}^{n} \alpha_i v_i,\, v_2, \dots, v_n\right)
            = \omega(0, v_2, \dots, v_n)
            = \omega(0 \cdot 0, v_2, \dots, v_n)
            = 0 \cdot \omega(0, v_2, \dots, v_n) = 0        
        \]
        \\

        \textbf{\emph{1}} (полилинейность по первому аргументу, остальные аналогично): \\
        Хотим $\omega(\alpha v + \beta г, \dots) = \alpha \omega(v, \dots) + \beta \omega(u, \dots)$
        Для этого надо $\omega(v + u, \dots) = \omega(v, \dots) + \omega(u, v_2, \dots, v_n)$.

        \textbf{1.а)} Если $v, v_2, \dots, v_n$ и $u, v_2, \dots, v_n$ лин.\ зависимы,
        то $v + u,\, v_2, \dots, v_n$ тоже линейно зависимы, и равенство $0 = 0 + 0$ очевидно.

        \textbf{1.б)} НУО $v, v_2, \dots, v_n$ лин.\ независимы $\Rightarrow$ базис $V$,
        поэтому $u = \lambda v + \sum_{i=2}^{n} \alpha_i v_i$. Тогда:
        \begin{align*}
            \omega\!\left(v + u,\, v_2, \dots, v_n\right)
            &= \omega\!\left(v + \lambda v + \sum_{i=2}^{n}\alpha_i v_i,\, v_2, \dots, v_n\right) \\
            &= \omega\!\left((1+\lambda)v,\, v_2, \dots, v_n\right) \\
            &= (1+\lambda)\,\omega(v, v_2, \dots, v_n) \\
            &= \omega(v, v_2, \dots, v_n) + \omega\!\left(\lambda v + \sum_{i=2}^{n}\alpha_i v_i,\, v_2, \dots, v_n\right) \\
            &= \omega(v, v_2, \dots, v_n) + \omega(u, v_2, \dots, v_n) \\
        \end{align*}

        \textbf{\emph{2}}:
        \begin{align*}
            \omega(\dots, u+v, \dots, u+v, \dots)
            &= \omega(\dots, u, \dots, u, \dots)
            + \omega(\dots, u, \dots, v, \dots) \\
            &\quad + \omega(\dots, v, \dots, u, \dots)
            + \omega(\dots, v, \dots, v, \dots) \\
            &= 0 + \omega(\dots, u, \dots, v, \dots)
            + \omega(\dots, v, \dots, u, \dots) + 0
        \end{align*}
        Но левая часть равна нулю по пункту 3 (два одинаковых аргумента), поэтому:
        \[
            \omega(\dots, u, \dots, v, \dots) + \omega(\dots, v, \dots, u, \dots) = 0
        \]
        \[
            \implies \omega(\dots, u, \dots, v, \dots) = -\omega(\dots, v, \dots, u, \dots) 
        \]

        Таким образом $\omega \in (\bigwedge^n V)^*$. \\


        \textbf{\emph{3}} в \textcircled{$\Rightarrow$}: \\
        Пусть $\omega \neq 0$, тогда
        $\omega$ --- базис $(\bigwedge^n V)^*$.

        Пусть $\exists\, v_1, \dots, v_n$ лин.\ независимы $\Rightarrow$ базис $V$,
        поэтому $\{v_1 \wedge v_2 \wedge \dots \wedge v_n\}$ --- базис $\bigwedge^n V$. Тогда:
        \[
            (v_1 \wedge \dots \wedge v_n)^* = \lambda\omega
        \]
        \[
            \omega(v_1 \wedge \dots \wedge v_n)
            = \frac{1}{\lambda}(v_1 \wedge \dots \wedge v_n)^*(v_1 \wedge \dots \wedge v_n)
            = \frac{1}{\lambda} \neq 0 \qquad \square
        \]

    \end{proof}
\end{stmt}



\begin{remark}
    Таким образом, любой объём --- это элемент $(\bigwedge^n V)^*$.
    Как следствие, если $\omega_1$ и $\omega_2$ --- объёмы ($\neq 0$), то
    $\omega_1 = \lambda \omega_2$ для некоторого $\lambda \in \mathbb{R}$.
\end{remark}


\begin{stmt}
    Пусть $V$ --- евклидово пространство размерности $n$,
    $\mathbf{e} = (e_1, \dots, e_n)$ --- базис и $\mathbf{v} = (v_1, \dots, v_n)$ --- набор векторов,
    где $v_i = \sum_{j=1}^{n} c_{ji} e_j$, то есть $\mathbf{e} \cdot C_{\mathbf{e}\mathbf{v}} = \mathbf{v}$. Тогда:
    \[
        \omega(v_1, \dots, v_n) = \omega(e_1, \dots, e_n) \cdot \det C_{\mathbf{e}\mathbf{v}}
    \]

    \begin{proof}
        \emph{TODO} \\
    \end{proof}
\end{stmt}